\subsection{完备滤过模的线性紧性性质}

回忆一下，若 E 是 A-模，则 E 的仿射子集（或仿射线性簇）是任意空集或形如 \(a + M\) 的子集 F，其中 \(a \in E\)，而 M 是 E 的子模，称为 F 的方向（《代数》，第 II 章，§ 9，第 1、3 节）。

命题 7. 设 A 为滤过环，E 为滤过 A-模，\((\mathbf{E}_{n})\) 为 E 的滤过；假设 \(E_{0} = E\)，各 \(E_{n}\) 都是 E 的子模，各 A-模 \(E/E_{n}\) 都是 Artin 的，并且拓扑群 E 是 Hausdorff 且完备的。则 E 中非空闭仿射子集的任意递减序列的交非空。

	我们在第 6 节已经看到，由于 E 是 Hausdorff 且完备的，它可识别为 \(\tilde{\mathbf{E}} = \varprojlim \mathrm{E} / \mathrm{E}_n\)。设 \((\mathbf{W}_p)\) 是 E 中非空闭仿射子集的递减序列，并对所有 \(n \geqslant 0\) 令 \(\mathbf{W}_{p,n}\) 为 \(\mathbf{W}_p\) 在 \(\mathrm{E} / \mathrm{E}_n\) 中的典范像；我们将构造一个序列 \(x = (x_n) \in \tilde{\mathbf{E}}\)，使得 \(x_n \in \mathbf{W}_{p,n}\) 对所有 \(p\)、\(n\) 成立；于是 \(x \in \mathbf{W}_p + \mathbf{E}_n\) 对所有 \(p\)、\(n\) 成立，而由于 \(\mathbf{W}_p\) 是闭的，便有 \(x \in \mathbf{W}_p\) 对所有 \(p\) 成立（第 5 节），从而命题得证。

	由于 \(E/E_{0}=0\)，取 \(x_{0}=0\)。假设已对 \(x_{i}\)（\(0\leqslant i\leqslant n-1\)）作了定义，令 \(W_{p,n}^{\prime}\) 为 \(W_{p,n}\) 中其典范像在 \(E/E_{n-1}\) 中等于 \(x_{n-1}\) 的元素所成的集合；由于 \(x_{n-1}\in W_{p,n-1}\)，且 \(W_{p,n-1}\) 是 \(W_{p,n}\) 的典范像，\(W_{p,n}^{\prime}\) 非空，显然是 \(E/E_{n}\) 的仿射子集；此外，序列 \((W_{p,n}^{\prime})_{p\in\mathbf{N}}\) 递减。由于 \(E/E_{n}\) 是 Artin 的，该序列最终恒定（否则，作为 \(E/E_{n}\) 的方向的 \(W_{p,n}^{\prime}\) 的子模序列会严格递减，这是荒谬的）。于是只需取 \(x_{n}\) 为 \(\bigcap_{p\in\mathbf{N}}W_{p,n}^{\prime}\) 中的一个元素，然后即可归纳地构造 \((x_{n})\)。

命题 8. 假设 A 和 E 满足命题 7 的假设。设 \((\mathbf{F}_{p})\) 是 E 的闭子模的递减序列，且 \(\bigcap_{p} F_{p} = 0\)。那么，对于 E 中 0 的每个邻域 V，都存在 p 使 \(F_{p} \subset V\)（换言之，滤子 \((\mathbf{F}_{p})\) 的基收敛于 0）。

	不妨设 V 是某个 \(E_{n}\)，于是 E/V 是 Artin 的。记 \(F_{p}^{\prime} = (F_{p} + V)/V\)；由于 \(F_{p}^{\prime}\) 构成 E/V 的递减子模序列，存在整数 j，使 \(F_{p}^{\prime} = F_{j}^{\prime}\) 对所有 \(p \geqslant j\) 成立。我们将看到 \(F_{j}^{\prime} = \{0\}\)，这就完成证明。取 \(x \in F_{j}^{\prime}\)，并令 \(W_{p}\) 为 \(F_{p}\) 中其在 E/V 中的像为 \(x (p \geqslant j)\) 的元素所成的集合；按 j 的定义，\(W_{p}\) 是 E 的非空闭仿射子集，且显然 \(W_{p+1} \subset W_{p}\)；

于是由命题 7 可知，存在元素 \(y\)，属于所有 \(W_p\)。由于 \(W_p \subset F_p\) 且 \(\bigcap_{p \in \mathbf{N}} F_p = \{0\}\)，\(y = 0\)；由于 \(x\) 是 \(y\) 在 \(E / V\) 中的典范像，故 \(x = 0\)（参见~习题 15 至 21）。

